% Pedro Kroger 2010

\documentclass{article}
\usepackage{multicol}
\usepackage{color}
\usepackage{graphicx}
\usepackage{url}
\usepackage{dirtree}
\usepackage[landscape,top=1cm, bottom=2cm, left=2cm, right=2cm]{geometry}
\usepackage[T1]{fontenc}
\usepackage[scaled=0.85]{beramono}
%\usepackage{bera}
\usepackage{fourier}
\usepackage{tikz}

\setlength\columnseprule{.4pt}
\setlength\parindent{0pt}

\definecolor{pythonblue}{RGB}{54,112,160}

\renewcommand{\columnseprulecolor}{\color{pythonblue}}

\pagestyle{empty}

\definecolor{shade}{HTML}{D4D7FE}
%\renewcommand*\DTstyle{\seriesdefault\tiny}

\newcommand{\py}[2]{
  \texttt{#1}(\textit{#2})\\
}

\newcommand{\pyc}[1]{
  \texttt{#1}\\
}

\newcommand{\header}[1]{
  \begin{tikzpicture}
    \node [fill=pythonblue!20,rounded corners=5pt]
    {\parbox{.95\linewidth}{\large \textcolor{pythonblue}{\sf \textbf{\hspace*{1ex}\raisebox{-15pt}{#1}}}\vspace*{1ex}}
    };
  \end{tikzpicture}
  \par
}

\begin{document}

\begin{minipage}[b]{0.1\linewidth}
\includegraphics[scale=.64]{python-logo-notext}
\end{minipage}
\hfill
\begin{minipage}[b]{0.9\linewidth}
{\Huge \textcolor{pythonblue}{\textsf{\textbf{Python 3.1 Quick Reference}}}} \\
{\large available at \url{http://pedrokroger.net/python-quick-reference.html}} \\
\vspace{.8em}
\end{minipage}

\hrulefill

\begin{multicols*}{4}

\header{Built-in functions}

\py{abs}{x}
\py{all}{iter}
\py{any}{iter}
\py{ascii}{obj}
\py{bin}{x}
\py{bool}{[x]}
\py{bytearray}{[arg,\,[encoding,\,[errors]]]}
\py{bytes}{[arg,\,[encoding,\,[errors]]]}
\py{chr}{i}
\py{classmethod}{fn}
\py{compile}{source,\,filename,\,mode}
\py{complex}{real,\,[imag]}
\py{dict}{[arg]}
\py{dir}{[object]}
\py{divmod}{x,\,y}
\py{enumerate}{iter,\,start=0}
\py{eval}{source,\,[globals,\,[locals]]}
\py{filter}{fn,\,iter}
\py{float}{[x]}
\py{format}{alue,\,[format\_spec]}
\py{frozenset}{[iter]}
\py{getattr}{object,\,name,\,[default]}
\py{globals}{}
\py{hasattr}{object,\,name}
\py{hash}{object}
\py{help}{[object]}
\py{hex}{x}
\py{id}{object}
\py{input}{[prompt]}
\py{int}{[number | string,\,[base]]}
\py{isinstance}{object,\,classinfo}
\py{issubclass}{class,\,classinfo}
\py{iter}{object,\,sentinel}
\py{len}{object}
\py{list}{[iter]}
\py{locals}{}
\py{map}{func,\,iter,\,\ldots}
\py{max}{iter,\,[args...],\,*,\,[key]}
\py{memoryview}{obj}
\py{min}{iter,\,[args...],\,*,\,[key]}
\py{next}{iterator,\,[default]}
\py{object}{}
\py{oct}{number}
\py{open}{file, **keys}
\py{ord}{c}
\py{pow}{x,\,y,\,[z]}
\py{print}{[obj,\,\ldots,\,sep,\,end,\,file]}
\py{property}{**keys}
\py{range}{[start,] stop,\,[step]}
\py{repr}{object}
\py{reversed}{sequence}
\py{round}{x,\,[n]}
\py{set}{iter}
\py{setattr}{object,\,name,\,value}
\py{slice}{[start],\,stop,\,[step]}
\py{sorted}{iter,\,[key],\,[reverse]}
\py{staticmethod}{unction}
\py{str}{[object,\,[encoding,\,[errors]]]}
\py{sum}{terable,\,[start]}
\py{super}{[type,\,[object-or-type]]}
\py{tuple}{[iter]}
\py{type}{object}
\py{type}{name,\,bases,\,dict}
\py{vars}{[object]}
\py{zip}{*iters}

\header{Built-in constants}
\texttt{False}, \texttt{True}, \texttt{None}\\
\texttt{Ellipsis}, \texttt{NotImplemented}

\header{Boolean operations}

\begin{tabular}{ll}
\texttt{not x}      & \\ 
\texttt{x or y}     & \\ 
\texttt{x and y}    & \\ 
\end{tabular}

\header{Comparisons}

\begin{tabular}{ll}
\texttt{x < y}  & \\ 
\texttt{x <= y} & \\ 
\texttt{x > y}  & \\ 
\texttt{x >= y} & \\ 
\texttt{x == y} & \\ 
\texttt{x != y} & \\ 
\texttt{x is y} & \\ 
\texttt{x is not y} & \\ 
\texttt{x < y < z}  & \\ 
\end{tabular}

\header{Numeric operations}

\begin{tabular}{ll}
\texttt{x + y}   & sum \\
\texttt{x - y}   & subtraction \\
\texttt{x * y}   & multiplication \\
\texttt{x / y}   & quotient \\
\texttt{x // y}  & floored quotient \\
\texttt{x \% y}  & remainder \\
\texttt{-x}      & negation \\
\texttt{+x}      & identity \\
\texttt{x ** y}  & x to the power y \\
\end{tabular}

\header{Bit-string Operations}

\begin{tabular}{ll}
\texttt{x | y}   & or \\
\texttt{x \& y}  & and \\
\verb!x ^ y! & exclusive or \\
\verb!x << n! & bitwise left-shift \\
\verb!x >> n! & right-shift \\
\verb+~x+ & bitwise invert (integers)  \\ 
\end{tabular}

\header{Extended Assignment}

\begin{tabular}{lll}
\texttt{x += y}   & & \texttt{x /= y}   \\
\texttt{x \&= y}  & & \verb!x >>= y!    \\
\texttt{x -= y}   & & \texttt{x \%= y}  \\
\texttt{x |= y}   & & \verb!x <<= y!    \\
\texttt{x *= y}   & & \texttt{x **= y}  \\
\verb!x ^= y!    & & \texttt{x //= y} \\
\end{tabular}

\header{Sequence Assignment}

\texttt{w = [1, 2, 3, 4]}

\begin{tabular}{ll}
\texttt{a, *b = w} & \texttt{(1, [2, 3, 4])} \\
\texttt{a, *b, c = w} & \texttt{(1, [2, 3], 4)} \\
\end{tabular}

\header{Sequence Operations}

\begin{tabular}{ll}
\texttt{x in s} & membership \\
\texttt{x not in s} & membership \\
\texttt{s + t} & concatenation \\
\texttt{s * n, n * s} & copy s \texttt{n} times \\
\texttt{s[n]} & nth item of s \\
\texttt{s[i:j]} & from i to j \\
\texttt{s[i:j:k]} & from i to j, step k \\
\end{tabular}

\header{Slice examples}

\texttt{s = ['a', 'b', 'c']}

\begin{tabular}{lll}
\texttt{s[0]} & \texttt{'a'}  & first \\
\texttt{s[-1]} & \texttt{'c'}  & last \\
\texttt{s[1:]} & \texttt{['b', 'c']}  & rest \\
\texttt{s[:-1]} & \texttt{['a', 'b']}  & butlast \\
\end{tabular}

\header{Float methods}

\py{float.as\_integer\_ratio}{}
\py{float.hex}{}
\py{float.fromhex}{s}
\header{List methods}

\py{append}{obj}
\py{extend}{iter}
\py{count}{value}
\py{index}{value,\,[start,\,[stop]]}
\py{insert}{pos, obj}
\py{pop}{[index]}
\py{remove}{value}
\py{reverse}{}
\py{sort}{key=None, reverse=False}

\header{String methods}

\py{capitalize}{}
\py{center}{width,\,[fillchar]}
\py{count}{sub,\,[start,\,[end]]}
\py{decode}{[encoding,\,[errors]]}
\py{encode}{[encoding[,errors]]}
\py{endswith}{suffix,\,[start,\,[end]]}
\py{expandtabs}{[tabsize]}
\py{find}{sub,\,[start,\,[end]]}
\py{format}{*args,\,**kwargs}
\py{index}{sub,\,[start,\,[end]]}
\py{isalnum}{}
\py{isalpha}{}
\py{isdigit}{}
\py{islower}{}
\py{isspace}{}
\py{istitle}{}
\py{isupper}{}
\py{join}{iterable}
\py{ljust}{width,\,[fillchar]}
\py{lower}{}
\py{lstrip}{[chars]}
\py{partition}{sep}
\py{replace}{old,\,new,\,[count]}
\py{rfind}{sub [,start [,end]]}
\py{rindex}{sub,\,[start,\,[end]]}
\py{rjust}{width,\,[fillchar]}
\py{rpartition}{sep}
\py{rsplit}{[sep [,maxsplit]]}
\py{rstrip}{[chars]}
\py{split}{[sep,\,[maxsplit]]}
\py{splitlines}{[keepends]}
\py{startswith}{prefix,\,[start,\,[end]]}
\py{strip}{[chars]}
\py{swapcase}{}
\py{title}{}
\py{translate}{table,\,[deletechars]}
\py{upper}{}
\py{zfill}{width}

\header{Set methods}

\py{add}{elem}
\py{clear}{}
\py{copy}{}
\py{difference}{other,\,\ldots}
\py{difference\_update}{other,\,\ldots}
\py{discard}{elem}
\py{intersection}{other,\,\ldots}
\py{intersection\_update}{other,\,\ldots}
\py{isdisjoint}{other}
\py{issubset}{other}
\py{issuperset}{other}
\py{pop}{}
\py{remove}{elem}
\py{symmetric\_difference}{other}
\py{symmetric\_difference\_update}{other}
\py{union}{other,\,\ldots}
\py{update}{other,\,\ldots}

\header{Dictionary methods}

\py{clear}{}
\py{copy}{}
\py{fromkeys}{seq,\,[value]}
\py{get}{key,\,[default]}
\py{items}{}
%\py{iteritems}{}
%\py{iterkeys}{}
%\py{itervalues}{}
\py{keys}{}
\py{popitem}{}
\py{pop}{key,\,[default]}
\py{setdefault}{key,\,[default]}
\py{update}{[other]}
\py{values}{}
%\py{viewitems}{}
%\py{viewkeys}{}
%\py{viewvalues}{}
\header{File methods}

\py{close}{}
\py{flush}{}
\py{fileno}{}
\py{isatty}{}
\py{next}{}
\py{read}{[size]}
\py{readline}{[size]}
\py{readlines}{[sizehint]}
\py{xreadlines}{}
\py{seek}{offset, [whence]}
\py{tell}{}
\py{truncate}{[size]}
\py{write}{str}
\py{writelines}{sequence}

\header{From future (python 2.6)}

\texttt{from \_\_future\_\_ import <f>}

\begin{tabular}{ll}
\texttt{absolute\_import}    & PEP 328\\
\texttt{division}            & PEP 238\\
\texttt{print\_function}     & PEP 3105\\
\texttt{unicode\_literals}   & PEP 3112\\
\end{tabular}

\header{Keywords (keyword.kwlist)}
\begin{verbatim}
False None True and as assert
break class continue def del in
elif else except finally for is
from global if import lambda
nonlocal not or pass raise
return try while with yield
\end{verbatim}

\header{Built-in exceptions}

\dirtree{%
.1 BaseException.
.2 SystemExit.
.2 KeyboardInterrupt.
.2 GeneratorExit.
.2 Exception.
.2 StopIteration.
}

\dirtree{%
.1 BaseException.
.2 AssertionError.
.2 AttributeError.
.2 BufferError.
.2 ArithmeticError.
.3 FloatingPointError.
.3 OverflowError.
.3 ZeroDivisionError.
.2 EnvironmentError.
.3 IOError.
.3 OSError.
.3 WindowsError (Windows).
.3 VMSError (VMS).
.2 EOFError.
.2 ImportError.
.2 LookupError.
.3 IndexError.
.3 KeyError.
.2 MemoryError.
.2 NameError.
.3 UnboundLocalError.
.2 ReferenceError.
.2 RuntimeError.
.3 NotImplementedError.
.2 SyntaxError.
.3 IndentationError.
.3 TabError.
.2 SystemError.
.2 TypeError.
.2 ValueError.
.3 UnicodeError.
.3 UnicodeDecodeError.
.3 UnicodeEncodeError.
.3 UnicodeTranslateError.
.2 Warning.
.2 DeprecationWarning.
.2 PendingDeprecationWarning.
.2 RuntimeWarning.
.2 SyntaxWarning.
.2 UserWarning.
.2 FutureWarning.
.2 ImportWarning.
.2 UnicodeWarning.
.2 BytesWarning.
}

\end{multicols*}
\end{document}

\newpage

\header{The string module}

\pyc{ascii\_letters}
\pyc{ascii\_lowercase}
\pyc{ascii\_uppercase}
\pyc{digits}
\pyc{hexdigits}
\pyc{octdigits}
\pyc{printable}
\pyc{punctuation}
\pyc{whitespace}
\py{format}{}
\py{capwords}{}
\py{maketrans}{}


%%% Local Variables: 
%%% mode: latex
%%% TeX-master: t
%%% End: 
